\documentclass[11pt,letterpaper]{article}
\usepackage[T1]{fontenc}
\usepackage[utf8]{inputenc}
\usepackage{lmodern}
\usepackage[margin=1in,headheight=14pt]{geometry}
\usepackage{amsmath,amssymb,amsthm,mathtools}
\usepackage{booktabs,tabularx,array,longtable}
\usepackage{microtype}
\usepackage{xcolor}
\usepackage{enumitem}
\usepackage{xurl}
\usepackage[colorlinks=true,linkcolor=blue!45!black,urlcolor=blue!45!black,citecolor=blue!45!black]{hyperref}
\hypersetup{pdftitle={Beyond Dirichlet GRH: Methods and New Transfers},pdfauthor={Research analysis prepared for Vladimir Reshetnikov},pdfsubject={Sharp torsion specialization, analytic families, and standard Galois cohomology}}
\usepackage{fancyhdr}
\pagestyle{fancy}
\fancyhf{}
\lhead{\small Beyond Dirichlet GRH}
\rhead{\small Methods and new transfers}
\cfoot{\thepage}
\setlength{\parskip}{0.32em}
\setlength{\parindent}{0pt}
\setlength{\emergencystretch}{2em}
\setlist{itemsep=0.2em,topsep=0.35em}
\newtheorem{theorem}{Theorem}[section]
\newtheorem{proposition}[theorem]{Proposition}
\newtheorem{lemma}[theorem]{Lemma}
\newtheorem{corollary}[theorem]{Corollary}
\theoremstyle{remark}
\newtheorem{remark}[theorem]{Remark}
\newcommand{\Sha}{\operatorname{Sha}}
\newcommand{\Q}{\mathbb Q}
\newcommand{\C}{\mathbb C}
\newcommand{\Z}{\mathbb Z}
\newcommand{\src}[1]{\texttt{#1}}
\newcolumntype{Y}{>{\raggedright\arraybackslash}X}

\title{\textbf{Beyond Dirichlet GRH}\\[0.4em]
\large Reusable proof mechanisms, a sharp torsion specialization theorem,\\
and an explicit comparison with standard Galois cohomology}
\author{Research analysis prepared for Vladimir Reshetnikov}
\date{25 November 2026 \quad (Pacific date)}

\begin{document}
\maketitle
\thispagestyle{empty}

\begin{abstract}
We inspect the exact published revision of \texttt{lsernau/grh-lean} and
compare all five earlier transfer reports supplied through ProveIt. The
exported theorems establish RH and GRH for Dirichlet $L$-functions; the
present purpose is to reuse their internal arguments beyond familiar
consequences of those conclusions.

Three concrete transfers are proved. First, for any finite family of
integral $p$-adic formal power series with total interior witness valuation
at most $B$, a common nonvanishing torsion specialization exists at depth
$H=h_*(r)+\lfloor B\rfloor$. This depth is independent of the number of
variables and of degrees, and is sharp for one exponent under the stated
radial restriction. A simultaneous first-$m$-powers version gives
$H=h_*(r/m)+\lfloor mB\rfloor$. Second, the repository's stronger
valuation-approximation theorem extends from polynomials to integral formal
series over $\C_p$ with exactly the same torsion-order bound, yielding
simultaneous determinantal and Fitting-ideal bounds. Third, for odd $\ell$
and positive Tate exponent, the direct limit of the source's odd cocycle
modules is identified with a precise eigenspace of the standard locally
trivial $H^1$ group, transferring its uniform finite-level length bound.

These are proved research deductions and completed comparison arguments.
They do not settle a named major open conjecture. We identify the exact
remaining obstacles for nonabelian GRH, arithmetic nonvanishing, Vandiver,
and $K$-theory, and distinguish proof, formal verification, and literature
priority throughout.
\end{abstract}

\vfill
\small
\textbf{Source revision:}
\path{72fead102f663d5c7a407964a13dc0cf9d977207}.\\
\textbf{Earlier-report snapshot:}
\path{VladimirReshetnikov/ProveIt}, revision
\path{a92443f00663b1b8b3c233d17df0ad49dc4f2b2c}.\\
\textbf{Verification:} source inspection, independent proof reviews, and
exact finite computations were performed. The new theorems in this report
have not been compiled in Lean. External verification of the published
RH/GRH source is cited, not represented as a local rebuild.
\normalsize
\newpage
\tableofcontents
\newpage

\section{Findings and the standard of evidence}
\label{sec:findings}

\subsection{Answer to the research question}

The repository contains reusable arguments with mathematical value outside
RH and Dirichlet GRH. The strongest immediately usable transfers in this
analysis are its integral torsion-selection machinery and its uniform
cocycle bounds. The determinant mechanism also provides a substantial
research program, but transplanting it to a new $L$-function requires new
arithmetic constructions.

The work below moves beyond a list of possible applications. It proves
specific theorems for analytic families, matrix ranks, Fitting ideals, and
standard Galois cohomology. It also sharpens the diagnosis of why several
tempting conjectural applications are not yet established. No proof of
general Hecke GRH, Vandiver, BSD, zero simplicity, or an equivalent new
major conjecture is claimed.

\begin{center}
\small
\begin{tabularx}{\textwidth}{>{\raggedright\arraybackslash}p{0.24\textwidth}YY}
\toprule
Result & What is established here & Status relative to earlier reports\\
\midrule
Sharp finite torsion test & Explicit $H=h_*(r)+\lfloor B\rfloor$, with
sharpness; finite power-orbit extension & Not treated in the five reports;
complete proof in Section~\ref{sec:sharp}\\
Analytic valuation transfer & Same upstream order bound for integral formal
series over completed $\C_p$ & New extraction and proof here;
Section~\ref{sec:series-specialization}\\
Matrices and finite complexes & Common torsion points preserve selected
ranks; controlled Fitting valuations and inverse denominators & Concrete
applications of the two torsion results\\
Odd cocycle comparison & Standard all-places local $H^1$ kernel with exact
complex-conjugation parity; inherited length bound & Supplies a comparison
not proved in the source; Section~\ref{sec:odd-sha-comparison}\\
Nonabelian GRH and other targets & Specific missing arithmetic hypotheses
and failure modes & Research targets, not additional conjectures solved\\
\bottomrule
\end{tabularx}
\end{center}

\subsection{What has been verified about the input}

The inspected repository revision is
\path{72fead102f663d5c7a407964a13dc0cf9d977207}, whose commit records
25 November 2026. Its two editions are pinned to Lean~4.35.0 and Mathlib
\path{56a005855ea17e0b0958f73ec801ae140b0a1679}.
The five exported statements are in \path{grh/src/Statements.lean} and
\path{grh/src/Main.lean}. They concern the Riemann zeta function and all
Dirichlet characters over $\Q$, including imprimitive characters.
They do not quantify over arbitrary Hecke or automorphic
$L$-functions~\cite{repo,shortverify}.

Independent checks of this exact shorter revision are reported by Andrew
Sutherland. The Lean community's statement concerns the original RH proof
and confirms both the intended statement and its verification. These
sources support accepting the RH/Dirichlet-GRH premise while retaining the
precise scope of the exported declarations~\cite{shortverify,community}.

Our own work inspected the central proof interfaces, their definitions, and
selected dependencies. A static import traversal confirms that the
torsion-selection and odd-cocycle results do not import the RH/GRH exports.
They are ingredients of the proof, not corollaries obtained by applying its
final zero-location theorem. Static import inspection is weaker than
kernel checking and is not described as such.

\subsection{Proof, verification, and novelty are separate questions}

Theorems proved in this report have full mathematical arguments, with two
independent reviews of the main transfer proofs. Their formalization is
outlined in the supplied artifacts; it has not been completed. The
included exact computations check examples and counterexamples, not
universal analytic or cohomological statements.

The five prior reports contain no development of the present torsion,
formal-series, or Fitting directions. This establishes a difference from
the supplied work, not world-wide novelty. Related torsion-rigidity and
simultaneous nonvanishing results already exist in work of Serban and
Kriz--Nordentoft~\cite{serban,kn}. The source itself identifies numerous
antecedents. We make no priority claim for the exact formulations here;
their value is that their hypotheses, proofs, quantitative bounds, and
connections to this source are explicit.

\section{What the five previous reports already accomplish}
\label{sec:prior}

All five ZIP archives at the supplied ProveIt location were downloaded and
their Git blob identities checked. The comparison below describes their
contents; it is not an independent certification of every classification
or numerical claim they contain. Full comparison notes and SHA-256 values
are included in the artifacts~\cite{incoming}.

\begin{center}
\small
\begin{tabularx}{\textwidth}{p{0.27\textwidth}Y}
\toprule
Archive & Content already developed\\
\midrule
\path{grh-lean-transfers.zip} & CM coefficient-field classifications,
weighted prime races, an $S_3$ splitting-weight detector, and an obstruction
to replacing its nonabelian trace by fixed congruence weights.\\
\path{grh-transfer.zip} & Arithmetic zero-sum criteria, abelian-field
examples, and a one-sided split-prime criterion for the zeta function of
an arbitrary finite Galois extension.\\
\path{grh_transfer_analysis.zip} & Source audit, sharp rank/index-dependent
determinant inequalities, basis-completion costs, and the discriminant
$-23$ cubic as a concrete non-Dirichlet target.\\
\path{grh_transfer_analysis_2026-11-25.zip} & Local twist/global height
interface, sparse auxiliary-family selection, positive nonabelian
detectors, and careful $K$-theory and witness-quantifier limitations.\\
\path{grh_transfers.zip} & Partial exponent-six class-group classification,
prime-race sign criteria, and analytic deductions for selected products.\\
\bottomrule
\end{tabularx}
\end{center}

In particular, the earlier determinant analysis already proves, under its
explicit hypotheses,
\[
 R\leq\frac{A+(B+2I)/\theta}{m},
\]
and its completion-cost variant. Generic one-sided zero-detection criteria,
including finite-Galois and $S_3$ cases, are also already present. Restating
these would not advance the requested program. We retain the determinant
architecture as context and direct the new work toward the analytic and
cohomological interfaces that those reports do not develop.

\section{The actual reusable mechanism in the RH proof}
\label{sec:mechanism}

\subsection{One determinant is bounded in opposite directions}

Write $d$ for target dimension, $n$ for the number of independent selected
images, $m$ for the number of theta lifts, $M$ for the main contribution,
$R$ for the prime remainder, and $C$ for a logarithmic squared covolume.
The source's three central inputs have the following shape:
\begin{align*}
 C&\leq n(M-mR+L),\qquad\delta\leq dI, &&\text{short forms and index},\\
 (1-4/b^2)d&\leq n\leq d, &&\text{enough forms},\\
 d(M-3L)&\leq C+2\delta, &&\text{long integral forms}.
\end{align*}
Here $L=fa\log^4b$, $I=10\log(1+2b^2)fa\log^3b$, and the source has
$M\geq0$. Comparison cancels the main term and gives
\[
 nmR\leq d(4L+2I).
\]
The dimension estimate makes $n\geq d/2$; the scale separation absorbs
the remaining factors, giving
\begin{equation}
 mR\leq fa\log^5 b.
 \label{eq:source-main}
\end{equation}
This is \src{New.star\_inequality}. Choosing a chain length comparable to
$\sqrt{x}$ is part of the arithmetic construction that ultimately yields
the critical exponent~\cite{mainineq,guide}.

The difficulty in a new application is to construct the same main term on
both sides with sufficiently small errors. An arbitrary new prime weight
cannot be inserted in \eqref{eq:source-main} while leaving its proof
unchanged.

\subsection{The portable twisted-lattice inequality}

The theorem \src{New.TwistedLattice.covolume\_add\_index\_nonneg} is genuinely
generic. It concerns linear maps $\Theta_\xi$ from one complex vector
space to a family of inner-product spaces, integral lattices at the
$p$-adic identifications, and rational positive-definite Gram data at one
base twist. If $n$ chosen vectors are completed to a basis of dimension
$d$, it gives the schematic inequality
\begin{equation}
 0\leq C+2\delta+(d-n)(b_{\rm size}+b_{\rm denom})+b_{\rm volume}.
 \label{eq:portable-lattice}
\end{equation}
Its inputs include injectivity of the lifts, all-conjugate size bounds,
integral denominator-clearing data for the added vectors, and heights of
the determinant ratios. A favorable allowed twist may be chosen
\emph{separately at each $p$-adic identification}. The matrix that clears
the index at that identification must work uniformly across all its
allowed twists~\cite{twisted}.

This is why simultaneous twisting and denominator control matter. A
different favorable twist at each prime is useful only because the global
height comparison links the resulting determinant ratios. Neither the
existence of local twists alone nor the usual untwisted product formula
supplies all of \eqref{eq:portable-lattice}.

\subsection{Two distinctions that survive every attempted transfer}

\textbf{Character order and conductor norm.} The general twisting theorem
constructs a character of $p$-power order and a suitable auxiliary prime
set. It bounds the number of parameters and various normalized losses.
It does not bound the norms of the auxiliary primes or the global
conductor. Consequently its torsion-order bound cannot be advertised as a
least-conductor theorem, a density theorem, or a conductor-aspect
subconvexity estimate~\cite{general}.

\textbf{Determinants and individual forms.} A determinant controls a product
of eigenvalues. It cannot bound each eigenvalue: $\operatorname{diag}(T,T^{-1})$
has determinant one and unbounded norm as $T\to\infty$. The source's
spanning theorem controls a large family through a second moment. It does
not isolate an arbitrary prescribed eigenform. Additional arguments would
be necessary for applications to individual automorphic $L$-values,
Ramanujan-type estimates, or derivatives at zeros.

\section{A sharp finite torsion test for simultaneous nonvanishing}
\label{sec:sharp}

The repository controls valuations to within an arbitrarily small error. When
an application needs only nonvanishing, a much smaller explicit bound is
available. The following theorem is developed here from the same
specialization viewpoint. Its proof is independent of RH and GRH. The claim
is a proved mathematical deduction, with sharpness as specified below; a
claim of priority over the literature is not made.

Fix a prime $p$, write $v(p)=1$ and $v(0)=+\infty$, and put
$\mathcal O=\mathcal O_{\mathbb C_p}$. For $r\in\mathbb Q_{>0}$ define
\begin{equation}
 h_*(r)=\min\left\{h\geq1:
       \frac1{p^{h-1}(p-1)}\leq r\right\}.
 \label{eq:hstar}
\end{equation}
In the theorem, different functions are allowed to have different reference
points. All valuations of reference values are finite.

\begin{theorem}[Explicit simultaneous nonvanishing]
\label{thm:sharp-nonvanishing}
Let $s\geq1$ and let $F_1,\ldots,F_q$ be a finite family in
$\mathcal O[[T_1,\ldots,T_s]]$. Suppose that points
$x_i\in\mathbb C_p^s$ satisfy
\[
 v(x_{ia})\geq r>0,\qquad
 B_0:=\sum_{i=1}^q v(F_i(x_i))\leq B<\infty.
\]
Set
\begin{equation}
 H=h_*(r)+\lfloor B\rfloor.
 \label{eq:linear-H}
\end{equation}
Then there are roots of unity $\zeta_1,\ldots,\zeta_s$, each of exact
order $p^{h_a}$ for some $h_*(r)\leq h_a\leq H$, such that
\[
 F_i(\zeta_1-1,\ldots,\zeta_s-1)\ne0\quad(1\leq i\leq q).
\]
Thus $\zeta_a^{p^H}=1$ and $0<v(\zeta_a-1)\leq r$. The depth $H$
depends only on $p,r,B$, not on the number of variables, the number of
functions separately from their aggregate budget, or polynomial degrees.
\end{theorem}

For an empty family the conclusion just means choosing any allowed torsion
tuple. The case $s=0$ reduces to nonzero constants and needs no roots of
unity. These harmless cases will be suppressed in the proof.

\subsection{Weighted Gauss valuation}

For a nonzero integral series $F=\sum_\alpha a_\alpha T^\alpha$ put
\[
 w_r(F)=\min_\alpha\bigl(v(a_\alpha)+r|\alpha|\bigr),
 \qquad w_r(0)=+\infty.
\]
The minimum is attained: coefficients are integral and $r|\alpha|$ tends
to infinity, so only finitely many degrees compete below any fixed finite
value. For integral series one has
\begin{align}
 w_r(F+G)&\geq\min(w_r(F),w_r(G)),\label{eq:gauss-add}\\
 w_r(FG)&=w_r(F)+w_r(G),\label{eq:gauss-mul}\\
 w_r(F)&\leq v(F(x))\qquad\text{if all }v(x_a)\geq r.
 \label{eq:gauss-eval}
\end{align}
The last inequality follows term by term and then by completeness.
For completeness, multiplicativity can be seen by choosing
$u\in\overline{\mathbb Q}_p$ with $v(u)=r$. The series $F(uT)$ has
coefficients tending to zero with total degree. After scaling $F(uT)$ and
$G(uT)$ to make their least coefficient valuations zero, their reductions
are nonzero polynomials over the residue field. Their product is nonzero,
since that polynomial ring is a domain. Hence the product also has least
coefficient valuation zero. Undoing the scaling proves
\eqref{eq:gauss-mul}.

\subsection{Integral division on a Cartesian grid}

The needed division statement does not require that $\mathcal O$ be
Noetherian or discretely valued.

\begin{lemma}[Integral division by a distinguished polynomial]
\label{lem:division}
Suppose $g(T)=T^d+a(T)\in\mathbb Z_p[T]$ is monic, $d\geq1$,
$\deg a<d$, and every coefficient of $a$ is divisible by $p$. Every
$F\in\mathcal O[[T_1,\ldots,T_s]]$ has a decomposition
\[
 F=g(T_1)Q+R,\qquad Q,R\in\mathcal O[[T_1,\ldots,T_s]],
 \qquad\deg_{T_1}R<d.
\]
Successive division in each variable gives
\begin{equation}
 F=\sum_{a=1}^s g(T_a)Q_a+R,
 \qquad\deg_{T_a}R<d\quad(1\leq a\leq s),
 \label{eq:multidivision}
\end{equation}
with all $Q_a$ integral and $R$ a polynomial in all variables.
\end{lemma}

\begin{proof}
Give bounded coefficient arrays the coefficient supremum norm. They form a
complete normed space over $\mathbb C_p$; its integral unit ball is the
set of integral formal series. Let $S_d$ discard the first $d$ powers of
$T_1$ and shift the remaining powers down by $d$. Let $M_a$ multiply by
$a(T_1)$. The operator $K=S_dM_a$ has norm at most $p^{-1}$, since every
coefficient of $a$ has valuation at least one. Therefore
\[
 Q=\sum_{n=0}^{\infty}(-K)^nS_dF
\]
converges in coefficient supremum norm to an integral series. It satisfies
$(1+K)Q=S_dF$. For $R=F-g(T_1)Q$ this identity says $S_dR=0$, so
$\deg_{T_1}R<d$. Division in $T_2$ of this remainder preserves its degree
bound in $T_1$, because the operators act only on $T_2$. Continue through
all variables to obtain \eqref{eq:multidivision}.
\end{proof}

\begin{lemma}[Vanishing on a finite product grid]
\label{lem:grid-ideal}
Under Lemma~\ref{lem:division}, suppose that $g$ has $d$ distinct roots
$A\subset\mathbb C_p$, all of positive valuation. If an integral formal
series $F$ vanishes at every point of $A^s$, then
\[
 F\in\bigl(g(T_1),\ldots,g(T_s)\bigr)
       \mathcal O[[T_1,\ldots,T_s]].
\]
\end{lemma}

\begin{proof}
All integral series in \eqref{eq:multidivision} converge on $A^s$. At a
grid point every $g(T_a)$ vanishes, so the remainder $R$ vanishes there.
It has degree less than $d$ in every variable. Repeated one-variable
interpolation implies $R=0$: fix all but the first coordinate, use $d$
distinct roots, and then repeat for the coefficient polynomials. No
integrality of interpolation denominators is needed; division has already
produced integral quotients before interpolation is used over the field.
\end{proof}

\subsection{Proof and sharpness of the depth bound}

\begin{proof}[Proof of Theorem~\ref{thm:sharp-nonvanishing}]
Write $h_*=h_*(r)$, $L=H-h_*+1=\lfloor B\rfloor+1$, and set
\begin{equation}
 g(T)=\prod_{h=h_*}^{H}\Phi_{p^h}(1+T),
 \label{eq:cyclotomic-grid-poly}
\end{equation}
where $\Phi_n$ is the $n$th cyclotomic polynomial. Modulo $p$,
$\Phi_{p^h}(1+T)=T^{\varphi(p^h)}$, so $g$ is monic with every
nonleading coefficient divisible by $p$. Its roots are precisely
$\zeta-1$ for primitive $p^h$-th roots $\zeta$ with $h_*\leq h\leq H$;
they are distinct and have valuations $\rho_h\leq r$.

For a linear factor $T-\alpha$ one has
$w_r(T-\alpha)=\min(r,v(\alpha))$. By multiplicativity,
\begin{equation}
 w_r(g)=\sum_{h=h_*}^{H}\varphi(p^h)\rho_h=L>B.
 \label{eq:grid-cost}
\end{equation}
Let $F=\prod_iF_i$. Again by multiplicativity and evaluation,
\[
 w_r(F)=\sum_iw_r(F_i)\leq\sum_iv(F_i(x_i))\leq B.
\]
If the desired tuple did not exist, at every point of the grid $A^s$ at
least one $F_i$ would vanish, so $F$ would vanish on the entire grid.
Lemma~\ref{lem:grid-ideal} would give $F=\sum_ag(T_a)Q_a$ with integral
$Q_a$. Since $w_r(Q_a)\geq0$, \eqref{eq:gauss-add} and
\eqref{eq:grid-cost} would imply $w_r(F)\geq L>B$, a contradiction.
\end{proof}

\begin{proposition}[Sharpness for one exponent]
\label{prop:sharpness}
The depth in \eqref{eq:linear-H} cannot be decreased uniformly, even for
one variable and one polynomial, when output roots must satisfy
$0<v(\zeta-1)\leq r$.
\end{proposition}

\begin{proof}
Put $k=\lfloor B\rfloor$. If $k\geq1$, take
\[
 F(T)=\prod_{h=h_*}^{h_*+k-1}\Phi_{p^h}(1+T),\qquad x=0.
\]
Then $v(x)=+\infty\geq r$, while $F(0)=p^k$ has valuation $k\leq B$.
Every permitted torsion point of order dividing $p^{h_*+k-1}$ is a root
of $F$. The next level, $h_*+k$, is therefore necessary. If $k=0$, the
level $h_*$ is already the least one meeting the radial requirement, so
it cannot be reduced either.
\end{proof}

This statement of sharpness concerns $m=1$ and the stated radial
restriction. It is not a claim that every individual family needs this
depth, nor that the first-$m$-powers bound below is optimal.

\subsection{One tuple for the first several powers}

\begin{corollary}[Simultaneous nonvanishing along a finite power orbit]
\label{cor:power-orbit}
Under the assumptions of Theorem~\ref{thm:sharp-nonvanishing}, let
$m\geq1$ and put
\begin{equation}
 H=h_*(r/m)+\lfloor mB\rfloor.
 \label{eq:power-H}
\end{equation}
There is a tuple with $\zeta_a^{p^H}=1$ and
$0<v(\zeta_a-1)\leq r/m$ such that
\[
 F_i(\zeta_1^j-1,\ldots,\zeta_s^j-1)\ne0
 \quad\text{for every }i\text{ and }1\leq j\leq m.
\]
\end{corollary}

\begin{proof}
For each $i,j,a$, choose $y_{ija}$ such that
\[
 (1+y_{ija})^j-1=x_{ia},\qquad v(y_{ija})\geq r/j\geq r/m.
\]
Such a choice exists. If $x_{ia}=0$, use $y_{ija}=0$. Otherwise all roots
of the monic integral polynomial $(1+Y)^j-(1+x_{ia})$ are integral, and
the sum of their valuations is $v(x_{ia})$. At least one root has
valuation at least $v(x_{ia})/j$.

The formal substitutions
\[
 G_{ij}(T)=F_i\bigl((1+T_1)^j-1,\ldots,(1+T_s)^j-1\bigr)
\]
are well defined integral series, since the substituted series have zero
constant term. Their values at $y_{ij}$ equal $F_i(x_i)$; convergence
also follows because all arguments have positive valuation. The aggregate
reference budget for the $mq$ functions $G_{ij}$ is at most $mB$.
Apply Theorem~\ref{thm:sharp-nonvanishing} with radius $r/m$ and budget
$mB$.
\end{proof}

\subsection{Ranks and exact complexes}

\begin{corollary}[Simultaneous rank preservation]
\label{cor:rank}
Let $A_b(T)$ be finitely many matrices over
$\mathcal O[[T_1,\ldots,T_s]]$. For each $b$ select an $r_b$-minor
$D_b(T)$ and an interior point $x_b$ with
$v(D_b(x_b))<\infty$. Put $B=\sum_bv(D_b(x_b))$, and assume all
coordinates of all $x_b$ have valuation at least $r>0$.
With $H$ as in \eqref{eq:power-H}, one tuple of torsion parameters gives
\[
 \operatorname{rank}_{\mathbb C_p} A_b(\zeta^j-1)\geq r_b
 \qquad(b,\ 1\leq j\leq m).
\]
If $r_b$ is the rank over the fraction field of the formal-series domain,
all these inequalities are equalities.
\end{corollary}

\begin{proof}
Apply Corollary~\ref{cor:power-orbit} to the selected determinants. For
the equality assertion, all larger minors are identically zero formal
series and hence vanish at every convergent specialization.
\end{proof}

For a finite complex of finite free modules, apply the corollary to one
generic-rank minor of each differential. Every resulting complex over
$\mathbb C_p$ has its generic ranks, so its cohomology dimensions are the
generic dimensions. In particular, a generically exact complex becomes
exact at every one of the first $m$ powers of the selected torsion tuple.
This conclusion concerns cohomology after passing to the field; it does
not assert that integral torsion in cohomology vanishes.

\subsection{Essential limitations and an exact example}

\paragraph{Integral coefficients are essential.}
The analytic function $\log(1+T)=\sum_{n\geq1}(-1)^{n+1}T^n/n$
vanishes at every $p$-power torsion point $\zeta-1$, but is nonzero.
Its coefficients do not have a uniform lower valuation bound. It therefore
does not satisfy the theorem. Conversely, any family with a known common
lower coefficient bound can be scaled to an integral family, but that
scaling changes the reference budget and must be included.

\paragraph{Existence is weaker than near preservation of valuations.}
The linear depth \eqref{eq:power-H} guarantees that the selected values
are finite, not that they are within $\varepsilon$ of their reference
valuations. The stronger conclusion is addressed separately in
Section~\ref{sec:series-specialization}. Neither theorem bounds the norm of
the conductor of a global arithmetic character without an additional
construction relating these torsion parameters to that character.

\paragraph{Exact finite illustration.}
Take $p=3$, $r=1$, $m=3$, and
\[
 A(X,Y)=\begin{pmatrix}1&X\\Y&XY+X\end{pmatrix},\qquad
 B(X,Y)=\begin{pmatrix}X-Y&0\\0&1\end{pmatrix}.
\]
Their determinants are $X$ and $X-Y$. Add the function
$F_3(X,Y)=(1+X)^3-1$. Reference points $(3,3)$, $(3,6)$, and $(3,3)$
give valuations $1,1,2$, so $B=4$. The guaranteed depth is
$h_*(1/3)+12=14$. The explicit tuple
$(\zeta_{27},\zeta_{27}^2)$ already works at depth $3$:
\begin{center}
\begin{tabular}{cccc}
\toprule
Power $j$ & $v(\det A)$ & $v(\det B)$ & $v(F_3)$\\
\midrule
1 & $1/18$ & $1/18$ & $1/6$\\
2 & $1/18$ & $1/18$ & $1/6$\\
3 & $1/6$ & $1/6$ & $1/2$\\
\bottomrule
\end{tabular}
\end{center}
Every entry is finite, so both matrices have rank two at all three powers.
The accompanying script checks these equalities by exact cyclotomic
resultants, and verifies eighteen finite sharpness examples. These examples
illustrate the theorems; the proofs above establish the universal assertions.

% Contribution for inclusion in the main report.
% Requires amsmath, amssymb, amsthm; environments theorem, proposition, corollary, remark.
% Author: source-transfer audit agent. No Lean build of the new extension was performed.

\section{From polynomial specialization to bounded analytic families}
\label{sec:series-specialization}

The cleanest reusable analytic statement in the repository is a simultaneous
specialization result for integral polynomials.  Its bound on the orders of the
roots of unity is chosen before the polynomials and is independent of their
degrees.  This order of quantifiers permits an extension to bounded formal
power series with the \emph{same} order bound.  We give the extension explicitly;
it is a deduction from a source theorem, not a claim that a named open conjecture
has thereby been settled.

Fix a prime $p$.  Write $\mathbb C_p$ for the completion of an algebraic closure
of $\mathbb Q_p$, $\mathcal O_{\mathbb C_p}=\{a:v(a)\geq0\}$, and
$v(p)=1$, $v(0)=+\infty$.  For a positive integer $h$ put
\[
  \rho_h=\frac1{\varphi(p^h)}=\frac1{p^{h-1}(p-1)}.
\]
A primitive $p^h$-th root of unity $\zeta$ satisfies
$v(\zeta-1)=\rho_h$.  Every series in
$\mathcal O_{\mathbb C_p}[[T_1,\ldots,T_s]]$ converges when all
$v(T_a)>0$: its coefficients are integral and terms of total degree $n$
have valuations tending uniformly to infinity as $n\to\infty$.

\paragraph{The exact polynomial input.}
The source theorem is
\path{New.Selmer.exists_rootsOfUnity_ellVal_eval_le}, in
\path{grh/src/New/Congruences/Selmer/GeneralTheorem.lean}~\cite{roots,general}.
Fix $s_0\in\mathbb N$, $m\geq1$, $r\in\mathbb Q_{>0}$, a finite
index set $I$, finite bounds $\beta_i\geq0$, and $\varepsilon>0$.
There is an integer $H$, depending only on these parameters, such that for
$s\leq s_0$, integral polynomials $P_i$ over $\overline{\mathbb Q}_p$,
and algebraic integral points $x_i$ satisfying
\[
 v(x_{ia})\geq r,\qquad v(P_i(x_i))\leq\beta_i,
\]
there is a tuple of roots of unity with
\begin{align}
 \zeta_a^{p^H}&=1,& v(\zeta_a-1)&\leq r/m,\label{eq:torsion-radius}\\
 v\bigl(P_i(\zeta_1^j-1,\ldots,\zeta_s^j-1)\bigr)
   &\leq v(P_i(x_i))+\varepsilon
 &&(i\in I, 1\leq j\leq m).\label{eq:polynomial-torsion}
\end{align}
The coefficient ring in this Lean theorem is the ring of integers in
$\overline{\mathbb Q}_p$, \emph{not} the completed ring
$\mathcal O_{\mathbb C_p}$.  The approximation in the following proof accounts
for this distinction.

The proof in \texttt{RootsOfUnity.lean} also supplies a sufficient explicit
criterion.  Put
\[
 B=\sum_{i\in I}\beta_i,\qquad
 \delta=\frac{\min(\varepsilon,1)}{s_0+2},\qquad
 h_* =\min\{h\geq1:\rho_h\leq r/m\}.
\]
It is sufficient to take $H\geq h_*$ with
\begin{equation}
 mB+1<\delta\sum_{h=h_*}^{H}\frac1{h+1}.
 \label{eq:harmonic-bound}
\end{equation}
Indeed, the source passes from $P_i$ to the $m$ polynomials obtained by the
substitutions $T_a\mapsto(1+T_a)^j-1$, uses reference points with valuations
at least $r/m$, and applies its one-exponent theorem with total budget $mB$.
The harmonic criterion is then precisely the criterion used in that theorem's
proof.  One convenient, deliberately coarse choice is
\begin{equation}
 H=\left\lceil(h_*+1)\exp\!\left(\frac{mB+1}{\delta}\right)\right\rceil.
 \label{eq:explicit-harmonic-depth}
\end{equation}
For this choice the harmonic sum is at least
$\log((H+2)/(h_*+1))$, making the inequality strict.
This estimate controls valuation losses; it should not be confused with the
much smaller bound available when only nonvanishing is required.

\begin{theorem}[The same torsion bound for integral formal power series]
\label{thm:series-same-bound}
Keep the parameters above and let $H\geq1$ be any bound having the polynomial
property \eqref{eq:torsion-radius}--\eqref{eq:polynomial-torsion}.
For $s\leq s_0$, suppose
\[
 F_i\in\mathcal O_{\mathbb C_p}[[T_1,\ldots,T_s]],\qquad
 x_i\in\mathbb C_p^s,\qquad
 v(x_{ia})\geq r,\qquad v(F_i(x_i))\leq\beta_i.
\]
Then there is a tuple $\zeta$ satisfying \eqref{eq:torsion-radius} and
\begin{equation}
 v\bigl(F_i(\zeta_1^j-1,\ldots,\zeta_s^j-1)\bigr)
 \leq v(F_i(x_i))+\varepsilon
 \qquad(i\in I, 1\leq j\leq m).
 \label{eq:series-torsion}
\end{equation}
The bound $H$ is unchanged: it has no new dependence on truncation degrees,
coefficients, reference points, or convergence rates of the series.
\end{theorem}

\begin{proof}
The empty family is immediate from the polynomial property, so assume $I$ is
nonempty.  Choose a real number
\[
 A>\max\bigl(r,\max_{i\in I}(\beta_i+\varepsilon)\bigr).
\]
Every nonzero coordinate of a point
\[
 t_j(\zeta)=(\zeta_1^j-1,\ldots,\zeta_s^j-1),\quad
 \zeta_a^{p^H}=1,\quad1\leq j\leq m,
\]
has valuation at least $\rho_H>0$.  If $\zeta_a^j=1$, the coordinate is zero
and its valuation is $+\infty$, so the same lower bound still holds.  Choose
$N$ so large that
\[
 (N+1)\min(r,\rho_H)>A.
\]
Let $P_i$ be the total-degree-$N$ truncation of $F_i$.  The ultrametric
inequality and completeness give
\begin{equation}
 v\bigl(F_i(x_i)-P_i(x_i)\bigr)>A,\qquad
 v\bigl(F_i(t_j(\zeta))-P_i(t_j(\zeta))\bigr)>A
 \label{eq:uniform-series-tail}
\end{equation}
for \emph{every} tuple of roots of order dividing $p^H$ and every indicated
$j$.  Thus truncation is uniform across the entire finite torsion grid,
before the source theorem chooses its tuple.

There are only finitely many coefficients in the $P_i$.  By density of
$\overline{\mathbb Q}_p$ in $\mathbb C_p$, approximate each coefficient by
an algebraic coefficient with error of valuation greater than $A$, obtaining
$Q_i\in\mathcal O_{\overline{\mathbb Q}_p}[T_1,\ldots,T_s]$.
Similarly choose algebraic $y_{ia}$ with
$v(y_{ia}-x_{ia})>A$.  Since $A>r>0$, all $y_{ia}$ remain integral and have
valuation at least $r$.

Integral polynomials are nonexpanding on integral polydiscs: by telescoping
monomials, if every coordinate changes by an amount of valuation greater
than $A$, then the polynomial value changes by an amount of valuation
greater than $A$.  Coefficient approximation has the same property, at
\emph{every} integral point.  Consequently
\[
 v\bigl(Q_i(y_i)-F_i(x_i)\bigr)>A,\qquad
 v\bigl(Q_i(t_j(\zeta))-F_i(t_j(\zeta))\bigr)>A
\]
on the same full torsion grid.  The first inequality and
$v(F_i(x_i))\leq\beta_i<A$ imply the exact equality
\[
 v(Q_i(y_i))=v(F_i(x_i))\leq\beta_i.
\]
Apply the polynomial theorem to $Q_i,y_i$ with the original parameters
and the original bound $H$.  It gives a tuple satisfying
\[
 v(Q_i(t_j(\zeta)))\leq v(F_i(x_i))+\varepsilon
 \leq\beta_i+\varepsilon<A.
\]
The second approximation inequality therefore gives
$v(F_i(t_j(\zeta)))=v(Q_i(t_j(\zeta)))$, which proves
\eqref{eq:series-torsion}.
\end{proof}

\begin{remark}[Why the hypotheses matter]
No restricted-series assumption is needed: integral coefficients and
strictly positive valuations of all arguments give convergence.  This would
be false for unrestricted substitution at the units $\zeta_a$ themselves;
the arguments here are $\zeta_a^j-1$.
Integrality of the coefficients is also substantive.  Without a uniform
lower bound on their valuations, no order bound can depend only on reference
values.  For a proposed $H$ and $x=p^R$ with $R\geq r$, the polynomial
\[
 \frac{(1+T)^{p^H}-1}{(1+x)^{p^H}-1}
\]
has value $1$ at $x$ but vanishes at every $\zeta-1$ with
$\zeta^{p^H}=1$.  Its denominator destroys the required integrality.
\end{remark}

\subsection{Simultaneous determinantal and Fitting profiles}
\label{subsec:fitting-profiles}

Let $A_a(T)$, for $a$ in a finite set, be matrices of possibly different
sizes $u_a\times v_a$ with entries in
$\mathcal O_{\mathbb C_p}[[T_1,\ldots,T_s]]$.  For
$1\leq k\leq\min(u_a,v_a)$ define
\[
 \mu_{a,k}(x)=\min_{\substack{I\subseteq\{1,\ldots,u_a\},\ |I|=k\\
                             J\subseteq\{1,\ldots,v_a\},\ |J|=k}}
                 v\bigl(\det A_a(x)_{I,J}\bigr).
\]
The minimum is $+\infty$ if every such minor vanishes.  If finite, it is
attained by at least one minor.  Evaluation on the open unit polydisc is a
ring homomorphism, so evaluating a formal determinant agrees with taking
the determinant after evaluation.

\begin{corollary}[One torsion tuple controls finitely many Fitting profiles]
\label{cor:fitting-torsion}
Choose a finite set $\mathcal K$ of matrix-size pairs $(a,k)$, reference
points $x_{a,k}$ with coordinate valuations at least $r>0$, and finite bounds
\[
 \mu_{a,k}(x_{a,k})\leq\beta_{a,k}.
\]
For every $m\geq1$ and $\varepsilon>0$, the bound $H$ in
Theorem~\ref{thm:series-same-bound}, with
$B=\sum_{(a,k)\in\mathcal K}\beta_{a,k}$, yields a single tuple $\zeta$
such that for all $(a,k)\in\mathcal K$ and all $1\leq j\leq m$,
\begin{equation}
 \mu_{a,k}(t_j(\zeta))
 \leq\mu_{a,k}(x_{a,k})+\varepsilon.
 \label{eq:fitting-profile-bound}
\end{equation}
Each selected Fitting ideal thus remains nonzero with a controlled valuation
at all these torsion specializations.  The reference point may differ for
each selected matrix and each selected minor size.
\end{corollary}

\begin{proof}
For each $(a,k)$ choose one formal $k$-minor $F_{a,k}$ attaining the finite
minimum at $x_{a,k}$.  Apply Theorem~\ref{thm:series-same-bound} to this
finite family.  At every output point, the minimum over all $k$-minors is
at most the valuation of the chosen witness minor.  This gives
\eqref{eq:fitting-profile-bound}.
\end{proof}

For the presented module
$M_a(x)=\operatorname{coker}(A_a(x):\mathcal O_{\mathbb C_p}^{v_a}
\to\mathcal O_{\mathbb C_p}^{u_a})$,
$\operatorname{Fitt}_{u_a-k}(M_a(x))$ is the ideal generated by the
$k$-minors.  A finite set of elements of a valuation ring generates the
principal ideal of its member of least valuation, so $\mu_{a,k}$ is exactly
the valuation of that Fitting ideal.  The statement is meaningful for these
finite presentations without any Noetherian assumption on
$\mathcal O_{\mathbb C_p}$.

Equivalently, at each specialization one can diagonalize a finite matrix
over the valuation ring by choosing a nonzero entry of smallest valuation
and eliminating its row and column, repeatedly.  If the resulting nonzero
diagonal entries have valuations
$\lambda_1\leq\cdots\leq\lambda_q$, then
\[
 \mu_{a,k}=\lambda_1+\cdots+\lambda_k\qquad(1\leq k\leq q).
\]
Thus the result controls cumulative invariant-factor valuations.  It does
\emph{not} give a separate upper bound for every $\lambda_k$: the lists
$(1,3)$ and $(0,4)$ have the second list's partial sums no larger, while its
last entry is larger.  Nor is $\mu$ a usual module length over the
nondiscrete ring $\mathcal O_{\mathbb C_p}$; the assertion is about
valuations of finitely generated determinantal ideals.

For a square matrix with a nonzero determinant at a reference point, the
case of its full determinant also bounds inverse denominators.  If
$v(\det A(x))\leq\beta$, then the selected torsion specializations have
nonzero determinant of valuation at most $\beta+\varepsilon$.
Their adjugate matrices remain integral, so every entry of their inverses
has valuation at least $-(\beta+\varepsilon)$.

\paragraph{Scope of the transfer.}
This is a concrete theorem for bounded analytic families and finite
presentations.  It can be applied to matrices presenting arithmetic modules
only after those matrices, their coefficient integrality, and their
reference minors have been supplied.  An Iwasawa-theoretic or Selmer-theoretic
interpretation is additional mathematical work.  In particular, the
corollary does not assert the existence of an arithmetic presentation that
has not been constructed, and it does not establish density statements or
bounded-conductor results for the auxiliary prime sets in the repository's
global twisting theorem.

\section{A completed comparison with standard Galois cohomology}
\label{sec:odd-sha-comparison}

The repository bounds an explicitly defined module of cocycles, but its
documentation does not identify that module with a standard
Tate--Shafarevich group. For odd coefficient primes and positive Tate
exponents, the following argument supplies such a comparison. Its ingredients
are standard cohomology; the contribution here is to apply them to the exact
definitions and quantifiers of this source. We make no claim that this
comparison theorem is new to the literature.

Fix a CM number field $E\subseteq\mathbb Q(\mu_N)$, an odd prime $\ell$, the
ring of integers $\mathcal O$ of a finite extension $F/\mathbb Q_\ell$, and
an integer $e>0$. Put
\[
 A_n=(\mathcal O/\ell^n\mathcal O)(-e),\qquad
 A=(F/\mathcal O)(-e).
\]
Thus $g$ acts by $\chi_\ell(g)^{-e}$, where $\chi_\ell$ is the cyclotomic
character. Let $c$ be complex conjugation, which normalizes $G_E$. At each
finite place $v$ choose a decomposition group $D_v$ and inertia group $I_v$;
for $v\mid\ell$ write $K_v=D_v\cap\ker\chi_\ell$.

Let $S_n$ be the $\mathcal O$-module of continuous cocycles $z:G_E\to A_n$
such that
\[
 z|_{I_v}=0\quad(v\nmid\ell),\qquad
 z|_{K_v}=0\quad(v\mid\ell),\qquad
 z(cgc)=-z(g)\quad(g\in G_E).
\]
After suppressing its one-point coordinate index, this is exactly
\texttt{New.Selmer.shaOdd} for $e>0$ in
\path{Lib/GaloisRep/Main/SelmerModules.lean}~\cite{oddsha}. Define
\[
 \Sha^1(E,A)=\ker\!\left(
 H^1(G_E,A)\longrightarrow\prod_v H^1(D_v,A)\right),
\]
with continuous cohomology and all places in the product. Infinite places
give no condition because $E$ is totally imaginary. The source tests all
closure places. Choosing one decomposition group at each place is
equivalent here: the tested inertia groups and cyclotomic kernels act
trivially on the coefficients, so the cocycle formula makes pointwise
vanishing on them invariant under conjugation.

\begin{theorem}[Exact comparison for the odd cocycle module]
\label{thm:odd-sha-comparison}
There is a canonical $\mathcal O$-linear isomorphism
\[
 \varinjlim_n S_n\ \simeq\
 \Sha^1\bigl(E,(F/\mathcal O)(-e)\bigr)^{c=(-1)^{e+1}}.
\]
The transition on coefficients is multiplication by $\ell$, identified with
the inclusions
$\ell^{-n}\mathcal O/\mathcal O\subseteq
 \ell^{-(n+1)}\mathcal O/\mathcal O$.
\end{theorem}

\begin{proof}
Set $(Tz)(g)=z(cgc)$. Since $\chi_\ell(cgc)=\chi_\ell(g)$, the involution
$T$ preserves cocycles. Every coboundary
\[
 (\delta a)(g)=\bigl(\chi_\ell(g)^{-e}-1\bigr)a
\]
is fixed by $T$. As $2$ is invertible, an odd coboundary is zero, and every
$T$-odd cohomology class has a unique odd representative, namely
$(z-Tz)/2$. The local conditions are preserved by $T$ because conjugation
permutes places and preserves the cyclotomic kernel.

For $v\mid\ell$, the image
$\Gamma_v=\chi_\ell(D_v)$ is an open subgroup of
$\mathbb Z_\ell^\times$. Its pro-$\ell$ subgroup $\Gamma_v^0\simeq
\mathbb Z_\ell$ has index prime to $\ell$. If $u$ is a topological
generator of $\Gamma_v^0$, then $u^{-e}-1\ne0$. Procyclic
cohomology and divisibility of $F/\mathcal O$ give
\[
 H^1(\Gamma_v^0,A)=A/(u^{-e}-1)A=0.
\]
Inflation--restriction through the finite prime-to-$\ell$ quotient yields
$H^1(\Gamma_v,A)=0$. A second inflation--restriction sequence, for
$1\to K_v\to D_v\to\Gamma_v\to1$, shows that
$H^1(D_v,A)\to H^1(K_v,A)$ is injective. Since $K_v$ acts trivially, all
coboundaries restrict to zero on it. Therefore the condition $z|_{K_v}=0$
is equivalent to local triviality of $[z]$ at $v$.

For $v\nmid\ell$, inertia acts trivially on $A$, while Frobenius acts by
$(Nv)^{-e}$. An unramified local class belongs to
\[
 H^1(D_v/I_v,A)=A/\bigl((Nv)^{-e}-1\bigr)A=0.
\]
The scalar is nonzero and $A$ is divisible. Conversely, a locally trivial
class has representatives zero on $I_v$, because coboundaries vanish there.
Thus the source's local conditions with coefficients in $A$ describe the
full local kernel $\Sha^1(E,A)$.

Finally, every continuous cocycle from compact $G_E$ to discrete $A$ has
finite image, hence takes values in
$\ell^{-n}\mathcal O/\mathcal O$ for some $n$. This identifies the direct
limit of the odd cocycle modules with the $T=-1$ part of that local kernel.
The natural action on cochains includes the coefficient action:
\[
 (c\cdot z)(g)=c\bigl(z(c^{-1}gc)\bigr)=(-1)^e(Tz)(g).
\]
Its eigenvalue is consequently $(-1)^{e+1}$, as claimed.
\end{proof}

The passage to divisible coefficients is essential. We have not identified
each finite $S_n$ with an ordinary finite-coefficient $\Sha^1$ group:
multiplication by the relevant nonzero local scalar need not be
surjective on $\mathcal O/\ell^n$. The argument also explicitly excludes
$\ell=2$, where division by two is unavailable, and $e=0$, where the
local scalar vanishes and the source uses a different local condition.

\begin{corollary}[A uniform standard cohomological bound]
\label{cor:odd-sha-bound}
Fix $E\subseteq\mathbb Q(\mu_N)$ as above and a nonnegative integer $t$. The source
theorem \texttt{New.exists\_length\_shaOdd\_le}~\cite{oddsha} supplies an integer
\[
 1\le X\le 2^{\,4(N+t+3)^{27}}h_E
\]
such that, for every odd prime $\ell$, every coefficient ring $\mathcal O$
allowed by that theorem, and every $0<e<t$,
\[
 \operatorname{length}_{\mathcal O}
 \Sha^1\bigl(E,(F/\mathcal O)(-e)\bigr)^{c=(-1)^{e+1}}
 \le v_{\mathcal O}(X).
\]
Here $v_{\mathcal O}$ is normalized at a uniformizer. In particular this
group is finite, is annihilated by $X$, and vanishes if $\ell\nmid X$.
Its order is at most $q^{v_{\mathcal O}(X)}$, where $q$ is the size of the
residue field of $\mathcal O$.
\end{corollary}

\begin{proof}
The source bound holds for every finite level $S_n$. Their transition maps
are injective on actual cocycles, so their lengths form a nondecreasing
sequence of integers bounded by $v_{\mathcal O}(X)$. The maximum is attained;
after that, every inclusion has quotient of length zero and is an
isomorphism. The direct limit therefore has the same finite bound. A module
of length $r$ over a discrete valuation ring is annihilated by the $r$th
power of its maximal ideal, proving the annihilator assertion and the
cardinality bound.
\end{proof}

The source's coefficient-ring hypothesis includes enough roots of unity of
order the exponent of $(\mathbb Z/N\mathbb Z)^\times$. It has not been
silently dropped. The argument proves eventual stabilization of the
finite-level cocycles but gives no effective bound on the stabilization
index from the length bound alone. The comparison above is a human-readable
proof; it has not been compiled as an additional Lean theorem.

\begin{corollary}[Integral odd cocycles vanish]
Under the same assumptions, a continuous cocycle $z:G_E\to\mathcal O(-e)$
with the same local and oddness conditions is zero.
\end{corollary}
\begin{proof}
Each reduction modulo $\ell^n$ lies in $S_n$ and is killed by $X$. Hence
$Xz(g)\in\bigcap_n\ell^n\mathcal O=0$ for every $g$. Since $X\ne0$ and
$\mathcal O$ is torsion-free, $z(g)=0$.
\end{proof}

\subsection{The exact limitation for Vandiver and \texorpdfstring{$K$}{K}-groups}

Alp\"oge's announcement tentatively suggests further $K$-group vanishing
and cyclicity consequences~\cite{alponge}. This comparison repairs a
definitional gap, but does not prove Vandiver's
conjecture. First, the bound allows nonzero groups at primes dividing $X$.
Second, $E,N,t$ are fixed before $\ell$ varies: a statement about almost all
coefficient primes for a fixed field does not control the diagonal choice
$E=\mathbb Q(\mu_\ell)$. Third, the parity is decisive. A class descended
from $\mathbb Q$ is fixed by the natural action of complex conjugation.
For even $e$, the source's odd cocycles instead correspond to the natural
minus eigenspace. At odd $\ell$ these eigenspaces intersect trivially.

The usual cohomological obstruction related to $K_{4m}(\mathbb Z)$ uses
$H^2(\mathbb Z[1/\ell],\mathbb Z_\ell(2m+1))$; its dual coefficient has
Tate exponent $-2m$. Thus the odd-cocycle bound by itself addresses the
opposite parity from this Vandiver obstruction. The repository's separate
even-class machinery would need additional analysis and arithmetic input.
No $K$-group vanishing or cyclicity claim is deduced here. The general
class-group and cohomological parity relationship is explicit in
Kurihara's 1992 paper, \emph{Some remarks on conjectures about cyclotomic
fields and $K$-groups of $\mathbb Z$}, especially its introduction and
Sections~1 and~3~\cite{kurihara}.


\section{Which further problems are now within reach?}
\label{sec:targets}

\subsection{A concrete conditional arithmetic application}

The matrix results give a useful exact specification for a next arithmetic
project. Suppose a finite family of arithmetic objects has finite free
complexes over
$\Lambda=\mathcal O_K[[T_1,\ldots,T_s]]$, where $K/\Q_p$ is finite.
Assume their degree-one cohomology is zero over $\operatorname{Frac}\Lambda$.
Suppose, in addition, that a proved control theorem identifies the
cohomology after character specialization with the arithmetic groups of
interest, including any local correction or exceptional factors.

Choose one generic-rank minor from every relevant differential, together
with interior nonzero witnesses. Include additional nonzero functions
whose vanishing would invalidate the control theorem. If the resulting
aggregate valuation budget is $B$, Corollary~\ref{cor:power-orbit} produces
a common character tuple of order dividing
\[
 p^{h_*(r/m)+\lfloor mB\rfloor}
\]
at whose first $m$ powers all the required ranks and nonvanishing factors
hold. The cohomology dimension formula then yields the desired vanishing
over $\C_p$. If integral sizes are required, the stronger
Fitting-profile theorem supplies the appropriate valuation control.

This is a proved \emph{conditional transfer criterion}. It does not assert
the needed arithmetic complexes, generic vanishing, interpolation
identities, or control theorem for an arbitrary elliptic curve or
automorphic representation. Supplying those is the remaining arithmetic
problem. The hypothesis of generic vanishing must never be disguised as a
consequence of selecting a torsion point.

\subsection{A first genuinely non-Dirichlet GRH benchmark}

Let $F=\Q(\alpha)$ with $\alpha^3-\alpha-1=0$. Its discriminant is $-23$.
The previous reports use its imaginary quadratic resolvent
$K=\Q(\sqrt{-23})$ and a nontrivial class-group character $\psi$ to write
\[
 \zeta_F(s)=\zeta(s)L_K(s,\psi).
\]
The remaining zero-location question is about a Hecke factor over $K$,
not a Dirichlet factor over $\Q$. This is an established benchmark from
the earlier reports, not a newly proposed target here~\cite{incoming}.

The principal-class selector in those reports gives an analytic criterion.
The missing work is an arithmetic realization of its divisibility gains
and the full short/enough/long comparison. Rational congruence and quadratic
splitting tests in the existing data do not automatically detect this
nonabelian splitting condition. The new torsion and cohomology results may
help with the denominator side, but they do not supply that detector.

A useful next deliverable would therefore be one concrete local
divisibility identity for the principal-class statistic, with a compatible
archimedean norm formula. If those formulas exhibit the same main term,
the generic twisted-lattice theorem becomes applicable. Another
equivalent prime-error reformulation would not fill this gap.

\subsection{Why several attractive conclusions remain unsupported}

\begin{center}
\small
\begin{tabularx}{\textwidth}{>{\raggedright\arraybackslash}p{0.24\textwidth}YY}
\toprule
Target & Relevant reusable material & Missing ingredient\\
\midrule
Nonabelian Hecke or Dedekind GRH & Twisted lattices, local twisting,
zero-detection criteria & A new prime detector and quantitative
short/enough/long estimates for it\\
Simultaneous arithmetic nonvanishing & Bounded torsion selection, rank and
Fitting control & Integral arithmetic presentations, nonzero reference
values, and proved specialization comparisons\\
Vandiver and $K$-group vanishing & Uniform cocycle bounds and the comparison
proved here & Correct parity sector, control of primes dividing $X$, and
field-uniform input for the diagonal cyclotomic family\\
Zero simplicity or pair correlation & No applicable direct output &
Derivative bounds or information correlating distinct zeros\\
Twin primes or strong prime correlations & Single weighted prime sums
appear in the source & New estimates for products such as
$\Lambda(n)\Lambda(n+2)$\\
Individual subconvexity or Ramanujan bounds & Determinant/second-moment
estimates & Isolation of individual forms and uniform conductor control\\
\bottomrule
\end{tabularx}
\end{center}

The table records what this analysis does and does not establish. It is
not a claim that these conjectures are logically independent of every
possible future development of the method.

\subsection{Literature position and a disciplined next step}

Serban's work studies torsion rigidity and separation for integral formal
power series; Kriz--Nordentoft already obtain arithmetic simultaneous
nonvanishing in substantial settings~\cite{serban,kn}. These precedents
rule out presenting ``torsion specialization of power series'' or
``simultaneous nonvanishing'' as new merely by name. The specific additions
here are an explicit sharp depth theorem with proof, the unchanged-bound
completion argument, its matrix consequences, and the exact application
of standard cohomology to the source's unusual cocycle definition.

For an immediate implementation project, formalize these precise
interfaces and instantiate them for one integral arithmetic presentation
already under good control. For a new major zero-location theorem, the
discriminant-$-23$ detector is a clear first test. The latter is a larger
arithmetic project; its central missing estimates are identified rather
than assumed.

\section{Reproducibility and verification boundaries}
\label{sec:reproduce}

The archive contains this manuscript and its component \LaTeX{} files,
the compiled PDF, exact computation scripts and outputs, source-audit and
provenance data, independent proof-review notes, and a pinned extraction
of the upstream modules needed for the principal reusable interfaces.
All copied upstream source files retain their original bytes and are
accompanied by the Apache-2.0 license, notice, and attribution. An adjusted
Lake configuration and an audit entry point are explicitly identified as
added files.

The extraction contains the full transitive local import closure of the
general torsion theorem, the odd-cocycle bound, and the twisted-lattice
theorem. Mathlib and its external packages are pinned by the original
manifest; their binaries are not bundled. The extraction is provided for
reproduction and was not built in this session.

\begin{center}
\small
\begin{tabularx}{\textwidth}{p{0.43\textwidth}Y}
\toprule
Check & Observed result\\
\midrule
Upstream revision and selected declaration locations & Exact revision
recorded; declarations found and hashed\\
Five incoming archives & All five match the recorded Git blob hashes\\
Local import closures & No RH/GRH endpoint in the extracted interfaces\\
Mathematical transfer proofs & Written in full; independently reviewed\\
Exact finite computations & Eighteen sharpness cases, three power
specializations, and matrix/valuation examples pass\\
PDF & Compiled, rendered, and visually checked\\
New Lean theorem compilation & Not performed; no such certificate claimed\\
Complete upstream kernel replay & Not performed here; independent
verification sources cited separately\\
Literature priority & Not established; related work identified\\
\bottomrule
\end{tabularx}
\end{center}

The exact computation script uses integer cyclotomic resultants. For
$\alpha\in\Q(\zeta_{p^h})^\times$ it computes
\[
 v_p(\alpha)=
 \frac{v_p\bigl(N_{\Q(\zeta_{p^h})/\Q}(\alpha)\bigr)}{\varphi(p^h)},
\]
using the unique prime above $p$ and its total ramification. Rational
expressions are evaluated by subtracting numerator and denominator
valuations. This avoids floating-point decisions about zeros.

\appendix
\section{Source map for the central claims}

All paths below are relative to \path{grh/src/} at the pinned source
revision. The JSON audit records exact declaration locations and file
hashes. Path names are evidence locators, not assertions of novelty.

\small
\begin{longtable}{p{0.49\textwidth}p{0.45\textwidth}}
\toprule
File & Role\\
\midrule
\endhead
\path{Statements.lean}, \path{Main.lean} & Exact RH and Dirichlet-GRH
endpoint statements\\
\path{New/Reduction/MainInequality.lean} &
\path{New.star_inequality}; main-term cancellation\\
\path{New/Reduction/Lattices/TwistedLattice.lean} &
\path{New.TwistedLattice.covolume_add_index_nonneg}; generic local-twist
determinant comparison\\
\path{New/Congruences/Selmer/RootsOfUnity.lean} &
\path{New.Selmer.exists_rootsOfUnity_ellVal_eval_le_add}; uniform
polynomial valuation selection\\
\path{New/Congruences/Selmer/GeneralTheorem.lean} &
\path{New.Selmer.exists_rootsOfUnity_ellVal_eval_le}; all first $m$ powers;
\path{exists_anticyclotomic_twist}\\
\path{Lib/GaloisRep/Main/SelmerModules.lean} &
\path{shaLocalSet}, \path{shaOdd}; exact local conditions and
precomposition parity\\
\path{Lib/Cyclotomic/OddSha.lean} &
\path{New.exists_length_shaOdd_le}; one integer bounding all finite
levels and coefficient primes\\
\bottomrule
\end{longtable}
\normalsize

\section{A short mathematical glossary}

\textbf{A $p$-adic valuation} measures divisibility by $p$: $v(p^ka)=k$
when $a$ is a unit. Values may be rational in an algebraic extension.
Large valuation means small $p$-adic absolute value. Setting $v(0)=\infty$
makes a finite valuation bound imply nonvanishing.

\textbf{Torsion specialization} here means substituting
$T_a=\zeta_a-1$, with $\zeta_a$ a root of unity of $p$-power order. It
does not mean substituting $T_a=\zeta_a$; the subtraction is essential for
convergence of arbitrary integral formal series.

\textbf{A Fitting ideal} of a finitely presented module is generated by
minors of a presentation matrix. Its valuation measures a cumulative
divisibility obstruction. For a square full-rank matrix its last such
quantity is the determinant valuation.

\textbf{A cocycle} is a function $z$ satisfying
$z(gh)=z(g)+g z(h)$. Cohomology identifies cocycles differing by
$g\mapsto ga-a$. The source's \path{shaOdd} consists of actual cocycles;
Section~\ref{sec:odd-sha-comparison} explains why a suitable odd subspace
has no nonzero coboundaries and how the local conditions compare.

\textbf{A normalized length} need not be an integer if it is divided by a
ramification index. Accordingly, a positive normalized Selmer length below
one does not by itself imply that the module is zero. The uniform
cohomological bound in Section~\ref{sec:odd-sha-comparison} uses the ordinary
integer-valued length over a discrete valuation ring and its corresponding
uniformizer valuation.

\begin{thebibliography}{99}
\bibitem{repo} L. Sernau and collaborators, \emph{GRH for Dirichlet
$L$-functions, in Lean 4}, repository revision
\path{72fead102f663d5c7a407964a13dc0cf9d977207}, 25 November 2026.
\url{https://github.com/lsernau/grh-lean/tree/72fead102f663d5c7a407964a13dc0cf9d977207}.

\bibitem{guide} Same repository, \path{grh/GUIDE.md}, \path{grh/INDEX.md},
and \path{grh/REFERENCES.md}, at the cited revision. The guide identifies
the three central theorems and explains the distinction between generic
background and specialized proof ingredients.

\bibitem{shortverify} A. Sutherland, \emph{Verification files for the
shorter Lean proof of RH and GRH}, exact revision \path{72fead10}.
\url{https://math.mit.edu/~drew/shortgrh/}. Consulted 25 November 2026
(Pacific date).

\bibitem{community} Lean FRO and Mathlib Initiative, \emph{Lean
verification of the Riemann Hypothesis}, 21 November 2026.
\url{https://leanprover-community.github.io/blog/posts/riemann-hypothesis/}.

\bibitem{alponge} L. Alp\"oge, announcement and methodological remarks,
November 2026. \url{https://alpo.ge/remark.pdf}. The discussion of chains
of theta lifts is on pp.~1, 3--4; the tentative $K$-group suggestion is
footnote~6 on p.~4.

\bibitem{incoming} V. Reshetnikov, ProveIt incoming research reports,
revision \path{a92443f00663b1b8b3c233d17df0ad49dc4f2b2c}.
\url{https://github.com/VladimirReshetnikov/ProveIt/tree/a92443f00663b1b8b3c233d17df0ad49dc4f2b2c/docs/incoming}.
Five archive identities and hashes are recorded in the accompanying
\path{artifacts/provenance.json}.

\bibitem{mainineq} Sernau repository, \path{grh/src/New/Reduction/MainInequality.lean},
\path{New.star_inequality}, at the revision in~\cite{repo}.

\bibitem{twisted} Sernau repository,
\path{grh/src/New/Reduction/Lattices/TwistedLattice.lean},
\path{New.TwistedLattice.covolume_add_index_nonneg}, at the same revision.

\bibitem{roots} Sernau repository,
\path{grh/src/New/Congruences/Selmer/RootsOfUnity.lean},
\path{New.Selmer.exists_rootsOfUnity_ellVal_eval_le_add}, at the same revision.

\bibitem{general} Sernau repository,
\path{grh/src/New/Congruences/Selmer/GeneralTheorem.lean},
\path{New.Selmer.exists_rootsOfUnity_ellVal_eval_le} and
\path{New.Selmer.exists_anticyclotomic_twist}, at the same revision.

\bibitem{oddsha} Sernau repository, \path{grh/src/Lib/Cyclotomic/OddSha.lean},
\path{New.exists_length_shaOdd_le}, and
\path{grh/src/Lib/GaloisRep/Main/SelmerModules.lean}, at the same revision.

\bibitem{serban} V. Serban, \emph{An infinitesimal $p$-adic multiplicative
Manin--Mumford conjecture}, Journal de Th\'eorie des Nombres de Bordeaux
\textbf{30} (2018), 393--408.
\url{https://www.numdam.org/item/JTNB_2018__30_2_393_0.pdf}.

\bibitem{kn} D. Kriz and A. C. Nordentoft, \emph{Horizontal $p$-adic
$L$-functions}, arXiv:2310.20678.
\url{https://arxiv.org/abs/2310.20678}.

\bibitem{kurihara} M. Kurihara, \emph{Some remarks on conjectures about
cyclotomic fields and $K$-groups of $\Z$}, Compositio Mathematica
\textbf{81} (1992), 223--236.
\url{https://www.numdam.org/item/CM_1992__81_2_223_0.pdf}.
\end{thebibliography}
\end{document}
